\section*{第 \ref{ch:g12:subduction-magmatism} 章——俯冲带的岩浆作用}

\begin{solution}{exo:g12:subduction-magmatism:1}
(a) \omterm{def:g12:subduction-magmatism:arc}{火山弧}；(b) \omterm{def:g12:subduction-magmatism:hydrous-melting}{含水熔融}（助熔熔融）；(c) \omterm{def:g12:subduction-magmatism:granodiorite}{花岗闪长岩}；
(d) \omterm{def:g12:subduction-magmatism:batholith}{岩基}；(e) \omterm{def:g12:subduction-magmatism:continental-growth}{大陆增生}。
\end{solution}

\begin{solution}{exo:g12:subduction-magmatism:2}
(1) 至少为 $24/0.25 = \qty{96}{km/h}$，即 $\num{96000}/3600 \approx
\qty{27}{m/s}$。(2) 爆炸速度至少快 $480/96 = 5$ 倍。
\end{solution}

\begin{solution}{exo:g12:subduction-magmatism:3}
(1) $838/1214 \approx 0.690$，约 \qty{69}{\percent}。(2) 在
上覆板块的一侧，即不下沉的那一侧。(3) 海沟与
前缘之间，\omterm{def:g11:subduction-zones:slab}{板片}尚不够深：还未到达
约一百千米的深度，\omterm{def:g11:subduction-zones:slab}{板片}中的水在那里释放到
足够热、含水即可熔融的地幔中。
\end{solution}

\begin{solution}{exo:g12:subduction-magmatism:4}
(1) $45/(45 + 15) = 0.75$：长石中有 \qty{75}{\percent} 为
斜长石。(2) 是的：该岩石为粗粒，石英含量超过
\qty{20}{\percent}，且其斜长石比例介于 65 与
\qty{90}{\percent} 之间。
\end{solution}

\begin{solution}{exo:g12:subduction-magmatism:5}
弧岩浆富含二氧化硅，因而动力黏度更大；同时它们携带
形成它们的水，因而富含气体。高动力黏度与
高气体含量都有利于爆炸式喷发。
\end{solution}

\begin{solution}{exo:g12:subduction-magmatism:6}
(1) $H = 150 \times \tan\ang{35} \approx 150 \times 0.700 \approx
\qty{105}{km}$。(2) 这是全球平均值。
\end{solution}

\begin{solution}{exo:g12:subduction-magmatism:7}
$d = H/\tan\theta$：$105/\tan\ang{20} \approx 105/0.364 \approx
\qty{288}{km}$（\ang{20} 时）；而 $105/\tan\ang{45} = \qty{105}{km}$（
\ang{45} 时）。\omterm{def:g11:subduction-zones:slab}{板片}越平缓，弧距海沟越远。
\end{solution}

\begin{solution}{exo:g12:subduction-magmatism:8}
(1) $T = -5.104 \times 3.16^2 + 132.9 \times 3.16 + 1120.7 \approx -51.0 +
420.0 + 1120.7 \approx \qty{1490}{\celsius}$。(2) 约高出 $1490 - 1000 =
\qty{490}{\celsius}$。(3) 干燥时不能：\qty{1300}{\celsius} 低于
\qty{1490}{\celsius}。湿润时可以：其温度高于 \qty{1000}{\celsius}。
\end{solution}

\begin{solution}{exo:g12:subduction-magmatism:9}
(1) 约为 $2.5/2.7 \approx 0.93$ 年，即约 340 天。(2) 此次崩塌
移动的是原已存在的旧岩石；弧则由地幔生成新岩石。只比较
这里只比较体积。
\end{solution}

\begin{solution}{exo:g12:subduction-magmatism:10}
(1) $2.7 \times \num{1000000} = \qty{2.7e6}{km^3}$。(2) $\num{9e9}/2.7
\approx \num{3.3e9}$ 年，即 33 亿年。(3) 该时间短于
地球年龄，因此弧完全可能形成了整个\omterm{def:g12:continental-crust:continental-crust}{大陆地壳}；但
\omterm{def:g11:subduction-zones:subduction}{俯冲}也会把地壳带回地幔，因此实际上大部分地壳已
被多次生成又破坏。
\end{solution}

\begin{solution}{exo:g12:subduction-magmatism:11}
(1) 浅源地震也发生在上覆板块中，以及两板块之间的
断层上；只有深源地震仅位于\omterm{def:g11:subduction-zones:slab}{板片}之中。
(2) $116 - 80 = \qty{36}{km}$，与 \qty{86.4}{km} 之比为：$\tan\theta = 36/86.4
\approx 0.417$，$\theta \approx \ang{22.6}$。
\end{solution}

\begin{solution}{exo:g12:subduction-magmatism:12}
在命题所述的直角三角形中，\omterm{def:g11:subduction-zones:slab}{板片}为斜边 $L$
$d$ 为与角 $\theta$ 相邻的直角边，故 $\cos\theta = d/L$，且 $L =
d/\cos\theta$。以日本为例，$L = 270/\cos\ang{20} \approx 270/0.940 \approx
\qty{287}{km}$。按每年 \qty{83}{mm} 计，$\num{2.87e8}/83 \approx \num{3.5e6}$
年，约 350 万年。
\end{solution}

\begin{solution}{exo:g12:subduction-magmatism:13}
(1) $V = \pi r^2 h = \pi \times 0.5^2 \times 8 \approx \qty{6.3}{km^3}$。
(2) 约为崩塌体积的 $6.3/2.5 \approx 2.5$ 倍。
(3) $P = \rho g z$：在 \qty{6}{km} 深处，$2830 \times 9.81 \times 6000 \approx
\qty{1.67e8}{Pa} = \qty{0.17}{GPa}$；在 \qty{14}{km} 深处，约
\qty{3.89e8}{Pa} = \qty{0.39}{GPa}。
\end{solution}

\begin{solution}{exo:g12:subduction-magmatism:14}
(1) 每年 $2.7 - 3.2 = \qty{-0.5}{km^3}$，即每百万年
损失 \qty{5e5}{km^3}。(2) $500 \times \num{5e5} = \qty{2.5e8}{km^3}$，
即占地壳的 $\num{2.5e8}/\num{9e9} \approx \qty{2.8}{\percent}$。(3) 两个
速率都只是估计值，其不确定度远大于两者之间相差的
每年 \qty{0.5}{km^3}：在误差范围内，二者达到平衡。
\end{solution}

\begin{solution}{exo:g12:subduction-magmatism:15}
(1) $\num{20000} \times 10 = \qty{2e5}{km^3}$。(2) $\num{2e5}/1 =
\num{200000}$ 年，即全球所有大陆弧合起来所需的时间。一
个弧只接收其中很小一部分岩浆，因此单个\omterm{def:g12:subduction-magmatism:batholith}{岩基}需要的时间
长得多：它由\omterm{def:g11:igneous-rocks:intrusion}{深成岩体}逐个累积而成，历经数百万年。(3) 至少
是其顶部以上 \qty{5}{km} 厚的岩石。
\end{solution}

\begin{solution}{pb:g12:subduction-magmatism:1}
\textbf{1.} $111.2 \times \cos\ang{39} \approx 111.2 \times 0.777 \approx
\qty{86.4}{km}$。

\textbf{2.} $(144.1 - 141.0) \times 86.4 \approx 3.1 \times 86.4 \approx
\qty{268}{km}$，约 \qty{270}{km}。

\textbf{3.} $(141.0 - 139.7) \times 86.4 \approx \qty{112}{km}$。

\textbf{4.} 位于海沟以西的上覆板块，即不下沉的那一侧，
太平洋板块\omterm{def:g11:subduction-zones:subduction}{俯冲}到它的下方。

\textbf{5.} 此处\omterm{def:g11:subduction-zones:slab}{板片}尚不够深，还不能把其中的水释放到
足以熔融的炽热地幔中。

\textbf{6.} $116 - 80 = \qty{36}{km}$。

\textbf{7.} 经度一度约 \qty{86.4}{km}。

\textbf{8.} $\tan\theta = 36/86.4 \approx 0.417$，故 $\theta \approx
\ang{22.6}$。

\textbf{9.} 该线每向东延伸 \qty{86.4}{km} 上升 \qty{36}{km}，因此它
在岩手以东 $98/0.417 \approx \qty{235}{km}$ 处到达地表，位于东经
$141.0 + 235/86.4 \approx \ang{143.7}$ 处，距海沟以西约 $(144.1 - 143.7) \times 86.4
\approx \qty{35}{km}$。

\textbf{10.} \omterm{def:g11:subduction-zones:slab}{板片}并非平面：它以近乎水平的姿态进入海沟，
随着下沉，弯曲程度逐渐增大。

\textbf{11.} $268 \times \tan\ang{20} \approx 268 \times 0.364 \approx
\qty{98}{km}$。

\textbf{12.} 两种方法结果一致，均为 \qty{98}{km}。

\textbf{13.} 两者均落在全球范围之内，略浅于
平均值 \qty{105}{km}。

\textbf{14.} 梅方位于岩手以西 \qty{112}{km}，因此其下方的\omterm{def:g11:subduction-zones:slab}{板片}
深度约为 $98 + 0.417 \times 112 \approx \qty{145}{km}$。

\textbf{15.} 在地幔中，位于地壳底部以下约
$98 - 35 = \qty{63}{km}$ 处。

\textbf{16.} 它来自海水，在地壳形成处附近的洋脊进入\omterm{def:g11:subduction-zones:slab}{板片}，
大部分保存在蛇纹石中。

\textbf{17.} $T = -5.104 \times 3.16^2 + 132.9 \times 3.16 + 1120.7 \approx
\qty{1490}{\celsius}$。

\textbf{18.} 由 $1490 - 1000 = 490$ 至 $1490 - 950 = \qty{540}{\celsius}$，
约 \qty{500}{\celsius}。

\textbf{19.} \omterm{def:g11:subduction-zones:slab}{板片}比周围地幔低数百度，即使有水也过冷而难以熔融；
真正熔融的是其上方的热地幔楔，它在被水浸润之后
才会熔化。

\textbf{20.} \omterm{def:g11:subduction-zones:slab}{板片}位于日本东北部的\textbf{火山锋下方约 \qty{100}{km}
处}（两种方法均为 \qty{98}{km}）。
\end{solution}
